% Generated from project outputs. Do not edit by hand.
\begin{table}[!htbp]
\centering
\begin{threeparttable}
\caption{Cohort cell sizes by subsample: 2021}
\label{tab:small-cells-2021}
\small
\setlength{\tabcolsep}{4pt}
\renewcommand{\arraystretch}{1.15}
\begin{tabularx}{\linewidth}{@{}Xrrrrr@{}}
\toprule
Sample & Cohorts & n \ensuremath{<} 30 & n \ensuremath{<} 20 & n \ensuremath{<} 10 & Missing SE \\
\midrule
General & 94 & 1 & 0 & 0 & 0 \\
Men & 94 & 10 & 2 & 0 & 0 \\
Women & 94 & 31 & 11 & 1 & 0 \\
Urban & 94 & 5 & 1 & 0 & 0 \\
Rural & 94 & 32 & 23 & 14 & 1 \\
Indigenous & 94 & 31 & 25 & 15 & 0 \\
Non-Indigenous & 94 & 4 & 1 & 0 & 0 \\
Bottom 99 percent & 94 & 2 & 0 & 0 & 0 \\
Bottom 90 percent & 94 & 9 & 2 & 0 & 0 \\
Bottom 50 percent & 94 & 33 & 24 & 15 & 0 \\
\bottomrule
\end{tabularx}
\begin{tablenotes}[flushleft]
\footnotesize
\item Cells are education-by-age cohorts; n is the unweighted number of surveyed respondents. Cutoffs are strict and nested. Missing SE counts missing labor-income mean standard errors. Subsamples overlap and must not be added together.
\end{tablenotes}
\end{threeparttable}
\end{table}
